In order to discuss mathematical concepts, a compact object is required. This is done by compressing information into a single mathematical expression. In order to discuss polynomials rings, we begin with defining least compressed case, monomials. Then, we consider a collection of monomnials, a polynomial. Finally, we can consider the compact object, polynomial rings.

\begin{definition}
A \textbf{monomial} in $\vectorComponents{x}$ is a product of the form
    \[
        \xalpha = x_1^{\alpha_1} \cdot x_2^{\alpha_2} \cdots x_n^{\alpha_n}
    \]
where all $ \alpha = \vectorComponents{\alpha}$ are nonnegative integers.
\end{definition}

As the root prefix "mono" suggests, there is only one of these terms $\xalpha$. They are independent and selc-contained. However, as we are creating the tools to generalize and deal with more abstract concepts, we also want to include the ability to capture any linear combination of these monomials.

\begin{definition}\label{def:polynomial}
A \textbf{polynomial} \textit{f} in $\vectorComponents{x}$ with coefficients, $\vectorComponents{a}$, in a field $k$ is a finite linear combination of monomials. 
    \[
        f = \sum_{\alpha} a_{\alpha} x^\alpha,\quad a_{\alpha} \in k
    \]
\end{definition}

With the polynomial definitions at hand, we can construct the final abstraction for our mathematics involving polynomials, the polynomial ring. This is defined as $\polyRing$. When leveraging this powerful abstraction, we capture $n$ dimensions of variables on any field $k$. This achieves the goal of compact analysis.  